% !TEX program = xelatex
% NOTE: 为了更好的 Unicode 支持，请使用 XeLaTeX（或 LuaLaTeX）编译。
\documentclass[11pt]{article}
\usepackage{fontspec}
\usepackage{xeCJK}
\setmainfont{CMU Serif}
\IfFontExistsTF{Noto Serif CJK SC}{\setCJKmainfont{Noto Serif CJK SC}}{%
  \IfFontExistsTF{Source Han Serif SC}{\setCJKmainfont{Source Han Serif SC}}{%
    \setCJKmainfont{FandolSong-Regular}% TeX Live fallback
  }%
}%
\usepackage[margin=1in]{geometry}
\usepackage{amsmath}
\usepackage{amssymb}
\usepackage{graphicx}
\usepackage{tikz-cd}
\usepackage{multicol}
\usepackage{float}
\usepackage{booktabs}
\usepackage{siunitx}

\sisetup{table-number-alignment=center}

\setlength{\parindent}{0pt}
\setlength{\parskip}{0.3\baselineskip}
\setlength{\textfloatsep}{5pt plus 1pt minus 1pt}
\setlength{\floatsep}{4pt plus 1pt minus 1pt}
\setlength{\intextsep}{4pt plus 1pt minus 1pt}
\setlength{\abovecaptionskip}{2pt}
\setlength{\belowcaptionskip}{1pt}
\setlength{\tabcolsep}{3pt}
\renewcommand{\arraystretch}{0.88}

\makeatletter
\renewcommand{\maketitle}{%
  \begin{center}
    {\LARGE \@title\par}
  \end{center}
  \vspace{-0.6\baselineskip}
}
\makeatother

\title{Idea报告V2}
\author{}
\date{}

\begin{document}

\maketitle

\section{研究背景}

传统推荐系统多采用 “召回–排序” 或 “匹配–打分” 范式，性能提升常依赖更大的 embedding表，但边际收益正在减弱。生成式推荐将任务改写为序列生成：先把 item 编码为离散Semantic ID（SID），再由自回归模型生成目标 item 的 SID。由于这一路线天然契合大语言模型训练范式，已成为当前研究热点。

MiniOneRec 提供了较完整且可复现的流程，覆盖 SID 构建与 SFT \cite{kong2025minionerec}；其推荐导向强化学习阶段采用 GRPO，而 GRPO 方法来源于 DeepSeekMath \cite{shao2024deepseekmath}。这一流程为我们提供了可复现的研究起点，也提示了一个有待验证的假设：在当前生成式推荐框架中，静态且以语义为主导的 tokenizer 可能已成为性能提升的关键瓶颈。

\section{复现结果}

目前在MiniOneRec的基础上，基于Qwen3-1.7B的base model，我们的复现结果如下

其中，传统序列推荐基线主要包括 GRU4Rec \cite{hidasi2016gru4rec}、Caser \cite{tang2018caser} 和 SASRec \cite{kang2018sasrec}。
生成式推荐基线主要包括 HSTU \cite{zhai2024hstu}、TIGER \cite{rajput2023tiger}、LCRec \cite{zheng2023lcrec}、BIGRec \cite{bao2023bigrec}、D3 \cite{bao2024d3}、S-DPO \cite{chen2024sdpo} 以及 MiniOneRec \cite{kong2025minionerec}。

\begin{table}[htbp]
  \centering
  \caption{Industrial 数据集上的性能对比\protect\footnotemark[1]}
  \small
  \begin{tabular}{l *{6}{S[table-format=1.4]} c}
    \toprule
    方法 & {HR@3} & {NDCG@3} & {HR@5} & {NDCG@5} & {HR@10} & {NDCG@10} & {超过指标数} \\
    \midrule
    复现结果 & 0.1004 & 0.0890 & 0.1198 & 0.0970 & 0.1513 & 0.1073 & - \\
    GRU4Rec & 0.0638 & 0.0542 & 0.0774 & 0.0598 & 0.0999 & 0.0669 & 6/6 \\
    Caser & 0.0618 & 0.0514 & 0.0717 & 0.0555 & 0.0942 & 0.0628 & 6/6 \\
    SASRec & 0.0790 & 0.0700 & 0.0909 & 0.0748 & 0.1088 & 0.0806 & 6/6 \\
    HSTU & 0.0927 & 0.0885 & 0.1037 & 0.0918 & 0.1163 & 0.0958 & 6/6 \\
    TIGER & 0.0852 & 0.0742 & 0.1010 & 0.0807 & 0.1321 & 0.0908 & 6/6 \\
    LCRec & 0.0915 & 0.0805 & 0.1057 & 0.0862 & 0.1332 & 0.0952 & 6/6 \\
    BIGRec & 0.0931 & 0.0841 & 0.1092 & 0.0907 & 0.1370 & 0.0997 & 6/6 \\
    D3 & 0.1024 & 0.0991 & 0.1213 & 0.0989 & 0.1500 & 0.1082 & 1/6 \\
    S-DPO & 0.1032 & 0.0906 & 0.1238 & 0.0991 & 0.1524 & 0.1082 & 0/6 \\
    MiniOneRec & 0.1143 & 0.1011 & 0.1321 & 0.1084 & 0.1586 & 0.1167 & 0/6 \\
    \bottomrule
  \end{tabular}
\end{table}

\begin{table}[htbp]
  \centering
  \caption{Office 数据集上的性能对比\protect\footnotemark[1]}
  \small
  \begin{tabular}{l *{6}{S[table-format=1.4]} c}
    \toprule
    方法 & {HR@3} & {NDCG@3} & {HR@5} & {NDCG@5} & {HR@10} & {NDCG@10} & {超过指标数} \\
    \midrule
    复现结果 & 0.1190 & 0.1055 & 0.1321 & 0.1109 & 0.1517 & 0.1171 & - \\
    GRU4Rec & 0.0629 & 0.0528 & 0.0789 & 0.0595 & 0.1019 & 0.0669 & 6/6 \\
    Caser & 0.0748 & 0.0615 & 0.0865 & 0.0664 & 0.1093 & 0.0737 & 6/6 \\
    SASRec & 0.0861 & 0.0769 & 0.0949 & 0.0805 & 0.1120 & 0.0858 & 6/6 \\
    HSTU & 0.1134 & 0.1031 & 0.1252 & 0.1079 & 0.1400 & 0.1126 & 6/6 \\
    TIGER & 0.0986 & 0.0852 & 0.1163 & 0.0960 & 0.1408 & 0.1002 & 6/6 \\
    LCRec & 0.0921 & 0.0807 & 0.1048 & 0.0859 & 0.1237 & 0.0920 & 6/6 \\
    BIGRec & 0.1069 & 0.0961 & 0.1204 & 0.1017 & 0.1434 & 0.1091 & 6/6 \\
    D3 & 0.1204 & 0.1055 & 0.1406 & 0.1139 & 0.1634 & 0.1213 & 0/6 \\
    S-DPO & 0.1169 & 0.1033 & 0.1356 & 0.1110 & 0.1587 & 0.1255 & 2/6 \\
    MiniOneRec & 0.1217 & 0.1088 & 0.1420 & 0.1172 & 0.1634 & 0.1242 & 0/6 \\
    \bottomrule
  \end{tabular}
\end{table}

\footnotetext[1]{“超过指标数”表示本复现结果优于该方法的指标数/总指标数（总指标数为 6，对应 HR@3、NDCG@3、HR@5、NDCG@5、HR@10、NDCG@10）；“优于”按严格大于（$>$）统计，不包含相等（$=$）。}

从结果来看，我目前的复现虽然还没有达到 MiniOneRec 论文中的最佳指标，但已经超过了大多数 baseline。我目前比较怀疑有两个主要原因导致和论文仍有差距：  
\begin{enumerate}
  \item 使用的 backbone 规模较小：本文使用 Qwen3-1.7B，而原论文使用 Qwen2.5-7B；
  \item 开源代码中的 alignment tasks 与论文正文描述并不完全一致，代码里实际使用的对齐任务数量更少。
\end{enumerate}

\section{MiniOneRec}

MiniOneRec 的主流程大致可以概括为下面四步： \cite{kong2025minionerec}：
\begin{enumerate}
  \item 用文本编码器将 title/description 映射到连续语义空间；
  \item 用 RQ-VAE \cite{lee2022rqvae} 量化为三层 SID（典型设定：$L=3, K=256$，其中 $L$ 表示 SID 长度（层数），$K$ 表示每层 code 的词表大小（codebook size））；
  \item 将 SID token 接入 LLM 词表进行 SFT；
  \item 在 constrained decoding（合法 SID 约束）下执行 GRPO 强化学习。
\end{enumerate}
同时，MiniOneRec 借鉴了 LCRec \cite{zheng2023lcrec} 中的方法，通过 alignment tasks 绑定自然语言空间与 SID 空间，保证 SID 与自然语言的对齐，避免 SID 完全变成对 LLM 来说没有语义的离散符号。


\section{Motivation}

MiniOneRec 中的 tokenizer 基于 RQ-VAE \cite{lee2022rqvae} 对语义文本向量进行离散量化。
这种方式目前有一些潜在问题：
\begin{enumerate}
  \item 目前的量化主要基于文本语义，协同信息没有直接进入 tokenizer；
  \item 不同 item 可能会被量化到相同 SID，出现 collision，会影响推荐结果；
  \item 同一个 code 在不同前缀下可能代表完全不同的语义方向，这会增加 LLM 的学习难度。
\end{enumerate}

\subsection{当前 SID 的初步诊断}

为了判断上述问题是否真实存在，而不仅仅停留在直觉层面，我们基于当前最优 Industrial 与 Office 结果，对现有 SID 结构和评测错误进行了诊断分析。诊断分为两层：一是对现有 item-SID 映射本身做统计，观察 collision 和 prefix 结构；二是结合 evaluate 结果中的预测 SID 与真实 target SID，分析错误是否真的集中发生在“相近前缀但未命中”的情形。

诊断结果表明，\textbf{collision 确实存在，但并不是当前最主要的问题}。在 Industrial 数据集上，现有 SID 的 collision rate 约为 $0.4341\%$；在 Office 数据集上约为 $0.4337\%$。这一结果说明不同 item 被映射到完全相同 SID 的现象并非不存在，但其绝对比例并不高，因此很难单独解释当前与论文结果之间的全部差距。

相比之下，\textbf{前缀语义不稳定与同层 code 的上下文歧义更值得关注}。例如，在 Industrial 数据集中，某些二级 code（如若干 $\texttt{<b\_*>}$）会出现在多达 16 个不同的一级前缀下，而三级 code $\texttt{<c\_0>}$ 甚至出现在 96 个不同的二级前缀上下文中；在 Office 数据集中，类似现象同样存在，例如某些二级 code 会出现在 31 个不同一级前缀下，而 $\texttt{<c\_0>}$ 出现在 60 个不同的二级前缀上下文中。这意味着同一个 code 的语义高度依赖其历史前缀，而不是全局稳定的。

更重要的是，这种 prefix ambiguity 已经直接体现在真实预测错误中。在 Industrial 数据集的 top-1 错误中，约有 $21.5\%$ 的错误样本与真实 target 共享相同的一级前缀，约有 $7.7\%$ 的错误样本甚至与真实 target 共享前两级前缀；在 Office 数据集中，这两个比例分别约为 $10.2\%$ 和 $1.9\%$。同时，在 beam 内部，Industrial 有约 $41.8\%$ 的样本至少出现过一个与 target 共享一级前缀的候选，约 $25.3\%$ 的样本至少出现过一个共享前两级前缀的候选；Office 上对应比例分别约为 $42.8\%$ 和 $23.5\%$。这说明模型并不是总在“完全错误的语义子树”上进行搜索，很多时候它已经进入了目标 item 附近的 prefix 子空间，但仍然无法在局部子树内部完成精确区分。

这一现象也可以从具体案例中看出。例如，在 Industrial 数据集中，真实目标 item ``X-Treme Tape ... Black'' 的 SID 为 $\texttt{<a\_202><b\_198><c\_80>}$，而模型 top-1 预测为 ``X-Treme Tape ... Clear''，对应 SID 为 $\texttt{<a\_202><b\_198><c\_218>}$；两者前两级完全相同，仅最后一级不同。类似地，真实目标 ``3D Solutech Real White PLA'' 与错误预测 ``3D Solutech Silver Metal PLA'' 也共享相同的前两级 SID 前缀。在 Office 数据集中也有类似现象，例如真实目标 ``BIC ... Blue'' 与错误预测 ``BIC ... Red'' 共享相同的前两级前缀，真实目标 ``EXPO ... Assorted Colors'' 与错误预测 ``EXPO ... Black'' 也落在相同的前缀子树中。

进一步地，我们又做了一个\textbf{train-only 的协同兼容性相关性诊断}。需要强调的是，这一分析\textbf{不是对“SID 构建缺少协同信息”进行严格因果量化}，而是检验如下现象是否存在：模型 top-1 预测的 item 在文本上已经很像真实 target，但从训练交互统计得到的协同兼容性来看，真实 target 其实更符合用户历史。具体做法是仅用训练集构建 item-item 共现/转移强度分数，并逐条比较测试样本中真实 target 与 top-1 预测之间的协同得分差。对于 top-1 预测 SID 对应多个 item 的碰撞场景，我们不再直接忽略，而是采用保守的 \textit{best-case} 口径：将该预测 SID 下协同分最高的候选视为模型最有利的比较对象。即便如此，在 Industrial 数据集上，仍有约 $9.8\%$ 的 top-1 错误样本满足“真实 target 的协同分高于预测 SID 下最强候选”；在 Office 数据集上，这一比例约为 $10.8\%$。如果只看同前缀错误，这种现象会更加明显：在 Industrial 的 same-$l_1$ 与 same-$l_2$ 错误中，这一比例分别约为 $19.0\%$ 和 $22.8\%$；在 Office 上分别约为 $21.7\%$ 和 $40.5\%$。

同时，这些“协同上应更偏向 target”的错误样本在文本空间中往往仍然高度相似：在保守口径下，Industrial 中这类错误的 target 与预测候选之间平均文本余弦相似度约为 $0.840$，Office 上约为 $0.823$；而在 same-$l_1$ / same-$l_2$ 错误子集内，平均文本相似度进一步上升到约 $0.926/0.982$（Industrial）和 $0.939/0.984$（Office）。这说明当前系统的一部分错误并不是“完全偏离语义方向”，而更像是：模型已经找到了文本上非常接近的候选，却没有充分利用用户交互协同关系把真正更合理的目标 item 排到前面。换句话说，\textbf{当前 SID/tokenizer 更擅长组织文本语义近邻，而对行为上应当区分的细粒度候选仍缺乏足够敏感性}。

因此，基于当前诊断结果，我们认为现阶段最值得优先解决的，并不是“完整 SID 碰撞太多”，而是：\textbf{现有 SID 的前缀结构尚未形成足够稳定且可解释的层级语义，同层 code 在不同上下文下被反复复用，导致模型虽然常常进入正确的语义子树，却难以在局部子树内部进一步精确区分目标 item。} 这也为后续围绕 prefix consistency、全局对齐量化以及更推荐原生的 tokenizer 设计提供了更直接的动机。


\section{Idea}

结合当前复现结果与问题诊断，本文当前阶段聚焦于\textbf{SID / Tokenizer 的改进}，暂不将强化学习阶段作为主要创新点。整体思路是：先提升物品连续表征的推荐原生性，再改进离散量化方式，使生成式推荐中的 SID 具备更强的协同感知能力、更稳定的前缀语义以及更低的碰撞风险。
\par
受 ReSID \cite{liang2026resid} 启发，本文考虑引入 FAMAE 模块。

\subsection{FAMAE（Field-Aware Masked Auto-Encoding）}

在 ReSID \cite{liang2026resid} 中，作者认为：在电商或内容推荐中，商品不仅仅是一段长文本描述，它天然带有\textbf{结构化特征字段}，比如：\texttt{Item\_ID}（商品ID）、\texttt{Category\_1}（一级类目：食品）、\texttt{Store}（店铺）等。

FAMAE 的核心思想是：item 不应该只被看成一段文本，它其实天然包含很多结构化信息，比如类目、属性、店铺等。  
如果 tokenizer 只看文本，那么它学到的更像“语义相似”，而不一定是“推荐里真正有用的相似”。

在原版 ReSID \cite{liang2026resid} 中，FAMAE 仅仅依赖数据集中天然的结构化特征，并将这些特征向量与位置编码进行简单的求和池化以形成输入给 Transformer 的词元表示。我们可以加以改进，把文本特征、协同特征和一些结构化字段一起作为输入，让 tokenizer 学到更适合推荐任务的 item 表示。

具体而言，我们引入了图结构特征（如使用 LightGCN \cite{he2020lightgcn} 提取的高阶协同向量），为每个物品 $i$ 构造多域特征集：
$$\mathbf x_i = (\mathbf x_i^{\text{txt}},\mathbf x_i^{\text{cat}},\mathbf x_i^{\text{attr}},\mathbf x_i^{\text{cf}},\mathbf x_i^{\text{graph}}).$$
\begin{itemize}
  \item $\mathbf x_i^{\text{txt}}$：文本编码器输出的连续语义向量。
  \item $\mathbf x_i^{\text{cf}}$：来自轻量图模型（如 LightGCN）提取的高阶协同表示。
  \item $\mathbf x_i^{\text{cat}},\mathbf x_i^{\text{attr}}$：类目层级与属性字段。
  \item $\mathbf x_i^{\text{graph}}$：拓扑结构摘要（如 item-item 共现图上的 kNN 邻居结构表示），用作显式拓扑监督。
\end{itemize}

原版 ReSID 对字段的融合方式比较简单，基本上是投影后直接相加。  
我这里想加一个轻量级门控，让模型自己决定：当前 item 更应该依赖文本信息，还是更应该依赖协同信息。我们首先为每个字段定义专属的投影层 $\phi_f$ 将其映射到统一的 $d$ 维空间 $\mathbf{v}_i^{(f)} = \phi_f(\mathbf{x}_i^{(f)})$，参考 PRISM \cite{fang2026prism} 中提出的自适应协同去噪（Adaptive Collaborative Denoising）机制。
具体而言，我们通过一个轻量级门控网络 $\phi_{gate}$ 计算协同特征的置信度向量 $\mathbf{g}_i \in (0, 1)^d$ ：
$$\mathbf{g}_i = \sigma(\phi_{gate}(\mathbf{v}_i^{\text{cf}}))$$
利用该门控，我们将高维连续的语义与协同特征进行动态加权，并与其他离散结构化特征叠加，从而得到最终的融合表征 $\mathbf{e}_i$：
$$\mathbf{e}_i = \underbrace{(1 - \mathbf{g}_i) \odot \mathbf{v}_i^{\text{txt}} + \mathbf{g}_i \odot \mathbf{v}_i^{\text{cf}}}_{\text{动态平衡连续语义与协同}} + \underbrace{\mathbf{v}_i^{\text{cat}} + \mathbf{v}_i^{\text{attr}} + \mathbf{v}_i^{\text{graph}}}_{\text{离散结构化特征}}$$
接下来，借鉴 ReSID \cite{liang2026resid} 的掩码预测机制构建序列输入。对于用户的历史交互序列样本，历史物品（位置 $1$ 到 $t-1$）直接使用上述完整的融合表征 $\mathbf e_i$。对目标 item，我随机遮掉一部分字段，让模型根据上下文去预测这些被遮掉的信息。我们根据概率分布 $p(\mathcal{M})$ 采样出一个需要被掩码的字段集合 $\mathcal{M}$。将这些字段\textbf{替换为该字段专属的掩码标记}后再次融合，得到带有掩码的目标表征 $\tilde{\mathbf e}_{i_t}$。

该序列被输入到 Transformer 中进行自注意力计算，输出包含用户历史交互上下文的目标隐藏状态 $\mathbf{h}_t$。随后，针对每个被掩码的字段 $f \in \mathcal{M}$，使用其专属的字段解码器得到预测结果 $\hat{\mathbf{x}}_{i_t}^{(f)}$。由于输入特征具有异构性，我们计算不同的损失函数：

离散字段（如类目、属性）计算交叉熵损失：
\[
\ell_f = -\log p_\theta \left( \mathbf{x}_{i_t}^{(f)} \mid \mathbf{h}_t \right).
\]
连续字段（如文本、CF 向量）计算均方误差（MSE）：
\[
\ell_f = \left\|\hat{\mathbf{x}}_{i_t}^{(f)} - \mathbf{x}_{i_t}^{(f)}\right\|_2^2 \quad
\]

FAMAE 的核心优化目标是最小化这些掩码字段的加权重构损失：
\[
\mathcal{L}_{\text{FAMAE}}(\theta) = \mathbb{E}_{(u, i_{1:t})} \, \mathbb{E}_{\mathcal{M} \sim p(\mathcal{M})} \left[ \sum_{f \in \mathcal{M}} w_f \, \ell_f \right].
\]
这里引入了模态平衡权重 $w_f$，其目的是在梯度回传时平衡不同模态的信息密度，防止高维文本的重建梯度过度主导网络。

此外，为了进一步将图结构显式注入连续表征中，我们在 FAMAE 训练时加入了一个基于 InfoNCE \cite{oord2018cpc} 的邻域对比损失 $\mathcal{L}_{\text{topo}}$。通过查表获取当前物品在图结构上的正样本邻居 $j^+$ 和负样本 $j^-$，拉近它们在隐空间中的距离：
\[
\mathcal{L}_{\text{topo}} = - \mathbb{E}_{i} \log \frac{\exp(\cos(\mathbf e_i,\mathbf e_{j^+})/\tau)}{\exp(\cos(\mathbf e_i,\mathbf e_{j^+})/\tau)+\sum_{j^-}\exp(\cos(\mathbf e_i,\mathbf e_{j^-})/\tau)}.
\]
最后，为了确保门控网络真正学习到数据中的流行度先验，我们引入一个辅助损失，强制门控的平均激活值 $\bar{g}_i$ 对齐物品在训练集中的归一化流行度 $p_i \in [0,1]$：
\[
\mathcal{L}_{pop} = \mathbb{E}_i\big[(\bar g_i - p_i)^2\big].
\]
其中，$\bar{g}_i$ 定义为该 $d$ 维门控向量在所有特征维度上的均值，即 $\bar{g}_i = \frac{1}{d}\sum_{k=1}^d g_{i,k}$；而 $p_i \in [0,1]$ 代表物品 $i$ 的真实流行度标量。在具体计算时，我们统计物品 $i$ 在训练集中的历史交互次数，并通过最大最小归一化（Min-Max Scaling）将其线性映射到 $[0,1]$ 区间。

最终，我们的结构化特征提取器总损失定义为以上三项之和：
\[
\mathcal{L}_{\text{total}} = \mathcal{L}_{\text{FAMAE}} + \lambda_1\mathcal{L}_{\text{topo}} + \lambda_2\mathcal{L}_{\text{pop}}.
\]



\subsection{基于树状结构的全局对齐量化与防碰撞正则化（GAOQ \& Anti-Collision Regularization）}
这一部分主要想解决两个核心问题：
\begin{enumerate}
  \item SID 的前缀应具有稳定且可解释的语义；
  \item 不同 item 不应过于容易被压缩到同一个 SID 上。
\end{enumerate}
在通过 FAMAE 提取出融合了多域信息的连续物品表征 $\mathbf{e}_i$ 后，我们需要将其离散化为 MiniOneRec 所需的固定长度（如 3 级）的语义 ID。在原版 MiniOneRec 及早期的生成式推荐模型（如 TIGER \cite{rajput2023tiger}）中，通常直接使用残差量化（RQ-VAE \cite{lee2022rqvae}）或层次 K-Means 。

然而，当我们在连续表征中显式注入了协同信号后，传统的量化方法可能会出现几个问题：
\begin{itemize}
    \item \textbf{前缀歧义}：传统残差量化在不同层级是独立分配索引的。这导致同一个数字代码在不同的父节点前缀下，代表完全不同的语义方向，这极大增加了大模型在自回归解码时的条件不确定性。
    \item \textbf{ID 碰撞与码分配偏置}：引入协同信号后，行为相似的物品在连续空间中会高度聚集。如果缺乏约束，大量不同的物品极易被分配到完全相同的 SID（即发生碰撞）；同时少数热门 Code 会霸占词表，导致长尾物品无法被生成。
\end{itemize}

为了解决上述问题，我们可以设计一个\textbf{减少碰撞的全局对齐量化模块}来替换传统的RQ-VAE，通过三个关键机制在数学上对 SID 的生成过程进行约束：

\subsubsection{基于匈牙利匹配的全局对齐（GAOQ）}
借鉴 ReSID \cite{liang2026resid} 中的 GAOQ（Globally Aligned Orthogonal Quantization）思想，我们希望同一层级的相同 Code 在不同前缀下尽量保持一致的语义方向。
具体而言，第一层采用标准全局聚类；在构建第 $\ell$ 层（$\ell \ge 2$）的量化树时，对于任意父节点 $p$，其下划分的子簇中心集合为 $\{\boldsymbol{\mu}_{p,k}\}$，父簇中心为 $\boldsymbol{\mu}_p$。我们首先计算\textbf{中心化残差方向}：
$$ \mathbf{d}_{p,k} = \boldsymbol{\mu}_{p,k} - \boldsymbol{\mu}_p $$
为保证全局对齐的稳定性并支持 mini-batch 训练，子簇中心 $\boldsymbol{\mu}_{p,k}$ 采用指数移动平均（EMA）方式更新。在每个 batch 中，将被路由到该子簇的连续向量 $\mathbf{e}_i$ 做均值聚合，并按动量 $\gamma$ 更新中心：
$$\mu_{p,k} \leftarrow \gamma\mu_{p,k} + (1 - \gamma) \cdot \text{mean}\{e_i \mid i \in \mathcal{B}_{p,k}\}$$
其中 $\mathcal{B}_{p,k}$ 表示当前 batch 中被分到父节点 $p$ 的第 $k$ 个子簇的样本集合。

随后，我们在该层全局维护一组数量为 $K$ 的全局参考锚点 $\{\mathbf{a}_{\ell,j}\}_{j=1}^K$。由于通常有 $B_p \le K$，对于每个父节点，我们通过求解匈牙利算法（Hungarian Algorithm），在保证单射的条件下最大化余弦相似度，将各个子簇强制对齐到统一的全局索引上：
$$ \pi_p = \arg\max_{\pi \in \Pi} \sum_{k=1}^{B_p} \cos(\mathbf{d}_{p,k}, \mathbf{a}_{\ell,\pi(k)}) $$
其中 $\Pi$ 代表所有从子簇集合 $\{1, \dots, B_p\}$ 到全局锚点集合 $\{1, \dots, K\}$ 的单射函数空间。最后，我们将第 $k$ 个子簇正式分配给 Code 索引 $j = \pi_p(k)$。让同一层的相同 code 在不同前缀下尽量保持一致的语义方向，从而提高 SID 序列的可预测性。

\subsubsection{基于边界的唯一性惩罚（Uniqueness Penalty）}
为了从连续空间的源头解决 ID 碰撞问题，我们可以参考 HiD-VAE \cite{fang2025hidvae} 引入了唯一性损失。在训练的 mini-batch 中，如果检测到两个不同的物品被分配了完全相同的 SID 序列，我们将它们记录为碰撞对集合 $\mathcal{P}_{\text{col}} = \{(i,j): i<j,\ \text{SID}(i)=\text{SID}(j)\}$。
对于这些碰撞对，我们在连续空间中施加一个基于 margin 的排斥损失。当两个不同 item 被量化到同一个 SID，且它们的余弦相似度超过安全边界 $m$，就对它们施加惩罚，从而降低后续继续发生碰撞的概率。
$$ \mathcal{L}_{\text{uniq}} = \sum_{(i,j)\in\mathcal{P}_{\text{col}}} \max\big(0,\ \cos(\mathbf{e}_i, \mathbf{e}_j) - m\big) $$

\subsubsection{码本多样性正则（Diversity Regularization）}
为缓解码分配偏置，并防止码本坍塌，我们可以借鉴 LETTER \cite{wang2024letter} 的动机，通过最大化每一层 Code 分配分布的熵，强制模型均匀地使用词表中的所有 Code。令 $p_{\ell,k}$ 为第 $\ell$ 层第 $k$ 个 Code 在全局的分配概率，多样性正则项（即分布熵）定义为：
$$ \mathcal{L}_{\text{ent}} = -\sum_{\ell=1}^L \sum_{k=1}^K p_{\ell,k} \log(p_{\ell,k} + \epsilon) $$

\subsubsection{综合量化目标}
最终，我们在进行向量离散化与词表构建时，不仅要最小化常规的向量量化与承诺损失（$\mathcal{L}_{\text{vq}}$），还要联合优化上述防碰撞与对齐约束。为保持符号闭合与阅读一致性，模块 B 的输入统一记为 FAMAE 输出的连续表征 $\mathbf{e}_i$。其总损失函数定义为：
$$ \mathcal{L}_{\text{quant}} = \mathcal{L}_{\text{vq}} + \lambda_{\text{uniq}}\mathcal{L}_{\text{uniq}} - \lambda_{\text{ent}}\mathcal{L}_{\text{ent}} + \lambda_{\text{disp}} \sum_{\ell=1}^L \sum_{1 \le j < k \le K} \cos^2(\mathbf{a}_{\ell,j}, \mathbf{a}_{\ell,k}) $$

\begin{itemize}
    \item \textbf{向量量化与承诺损失（$\mathcal{L}_{\text{vq}}$）}：在将连续物品表征 $\mathbf{e}_i$ 经匈牙利匹配分配给各层级的全局锚点后，我们得到其离散化向量 $\hat{\mathbf{e}}_i = \sum_{\ell=1}^L \mathbf{a}_{\ell, c_{i,\ell}}$。由于硬量化操作不可导，我们采用标准的量化承诺形式来提供梯度：
    $$ \mathcal{L}_{\text{vq}} = \|\text{sg}[\mathbf{e}_i] - \hat{\mathbf{e}}_i\|_2^2 + \beta\|\mathbf{e}_i - \text{sg}[\hat{\mathbf{e}}_i]\|_2^2 $$
    （其中 $\text{sg}[\cdot]$ 为停止梯度操作）。公式的第一项为字典学习损失，用于驱动离散锚点主动逼近连续空间的数据分布；第二项为承诺损失（Commitment Loss），约束底层编码器的输出不要偏离选定的锚点过远。
    
    \item \textbf{全局锚点最大散度正则（$\mathcal{L}_{\text{disp}}$）}：尽量让不同锚点彼此分散。具体做法就是惩罚不同锚点之间的余弦相似度平方和。直观上，这会鼓励 codebook 更均匀地被使用，减少少数热门 code 过度拥挤的情况。
\end{itemize}


通过上述设计，我们不仅为 MiniOneRec 提供了一套高度浓缩了协同知识的 SIDs，还预计能够显著降低前缀歧义与碰撞率，并降低 constrained decoding 场景下的解码不确定性。后续我们将重点通过 \textit{collision rate}、\textit{prefix conditional entropy} 以及下游推荐指标三类结果进行定量验证。

\subsection{潜在局限性}

本方案在实际训练与工程落地中仍可能面临以下挑战：

\subsubsection{静态 Tokenizer 与协同信号动态演化的问题}
本方案采用了两阶段训练范式，即先训练并冻结 Tokenizer，随后将固化的 SID 输入给 LLM 进行自回归生成与强化学习。Tokenizer 是静态的，但协同关系是动态的。

\subsubsection{全局图特征引入的泄露风险与对比公平性}
在 FAMAE的特征融合阶段，我们显式引入了 LightGCN 提取的高阶协同向量与图拓扑摘要。在严谨的序列推荐评估（如基于时间划分的 Leave-one-out）中，如果构建全图时无意间使用了验证集或测试集的交互边，将会导致严重的“时间泄露”，使评估指标虚高。此外，MiniOneRec 原版 Tokenizer 仅依赖文本特征，如果我们将融合了图特征的本方案直接与之对比，可能会被质疑性能的提升源于“特征信息的增加”而非 Tokenizer 架构本身的优越性。

\subsubsection{防碰撞惩罚的局部稀疏性与扩展能力}
我们参考 HiD-VAE 引入了基于 Margin 的唯一性排斥损失（Uniqueness Penalty）来解决 ID 碰撞问题。然而，该损失的触发机制高度依赖于在当前训练的 Mini-batch 内同时采样到被分配了相同 SID 的不同物品。在拥有数百万级物品的真实工业场景中，随机采样导致的 Batch 内硬碰撞概率极低。这种触发的稀疏性可能导致该惩罚项在实际训练中实际上没有作用，无法有效将全局潜空间中过度聚集的长尾物品推开。

\clearpage

\begin{thebibliography}{99}

\bibitem{hidasi2016gru4rec}
Bal\'azs Hidasi, Alexandros Karatzoglou, Linas Baltrunas, and Domonkos Tikk.
Session-based Recommendations with Recurrent Neural Networks.
In \textit{Proceedings of ICLR}, 2016.

\bibitem{tang2018caser}
Jiaxi Tang and Ke Wang.
Personalized Top-N Sequential Recommendation via Convolutional Sequence Embedding.
In \textit{Proceedings of WSDM}, 2018.

\bibitem{kang2018sasrec}
Wang-Cheng Kang and Julian McAuley.
Self-Attentive Sequential Recommendation.
In \textit{Proceedings of ICDM}, 2018.

\bibitem{he2020lightgcn}
Xiangnan He, Kuan Deng, Xiang Wang, Yan Li, Yongdong Zhang, and Meng Wang.
LightGCN: Simplifying and Powering Graph Convolution Network for Recommendation.
In \textit{Proceedings of SIGIR}, 2020.

\bibitem{lee2022rqvae}
Doyup Lee, Chiheon Kim, Saehoon Kim, Minsu Cho, and Wook-Shin Han.
Autoregressive Image Generation using Residual Quantization.
In \textit{Proceedings of CVPR}, 2022.

\bibitem{oord2018cpc}
Aaron van den Oord, Yazhe Li, and Oriol Vinyals.
Representation Learning with Contrastive Predictive Coding.
\textit{arXiv preprint arXiv:1807.03748}, 2018.

\bibitem{kuhn1955hungarian}
Harold W. Kuhn.
The Hungarian Method for the Assignment Problem.
\textit{Naval Research Logistics Quarterly}, 2(1--2):83--97, 1955.

\bibitem{shao2024deepseekmath}
Zhihong Shao, Peiyi Wang, Qihao Zhu, Runxin Xu, Junxiao Song, Mingchuan Zhang, Y. K. Li, Y. Wu, and Daya Guo.
DeepSeekMath: Pushing the Limits of Mathematical Reasoning in Open Language Models.
\textit{arXiv preprint arXiv:2402.03300}, 2024.

\bibitem{kong2025minionerec}
Xiaoyu Kong, Leheng Sheng, Junfei Tan, Yuxin Chen, Jiancan Wu, An Zhang, Xiang Wang, and Xiangnan He.
MiniOneRec: An Open-Source Framework for Scaling Generative Recommendation.
\textit{arXiv preprint arXiv:2510.24431}, 2025.

\bibitem{zheng2023lcrec}
Bowen Zheng, Yupeng Hou, Hongyu Lu, Yu Chen, Wayne Xin Zhao, Ming Chen, and Ji-Rong Wen.
Adapting Large Language Models by Integrating Collaborative Semantics for Recommendation.
\textit{arXiv preprint arXiv:2311.09049}, 2023.

\bibitem{rajput2023tiger}
Shashank Rajput, Nikhil Mehta, Anima Singh, Raghunandan Hulikal Keshavan, Trung Vu, Lukasz Heldt, Lichan Hong, Yi Tay, Vinh Q. Tran, Jonah Samost, Maciej Kula, Ed H. Chi, and Mahesh Sathiamoorthy.
Recommender Systems with Generative Retrieval.
In \textit{Advances in Neural Information Processing Systems 36 (NeurIPS 2023)}, 2023.

\bibitem{zhai2024hstu}
Jiaqi Zhai, Lucy Liao, Xing Liu, Yueming Wang, Rui Li, Xuan Cao, Leon Gao, Zhaojie Gong, Fangda Gu, Jiayuan He, Yinghai Lu, and Yu Shi.
Actions Speak Louder than Words: Trillion-Parameter Sequential Transducers for Generative Recommendations.
In \textit{Proceedings of ICML}, 2024.

\bibitem{bao2023bigrec}
Keqin Bao, Jizhi Zhang, Wenjie Wang, Yang Zhang, Zhengyi Yang, Yancheng Luo, Fuli Feng, Xiangnan He, and Qi Tian.
A Bi-Step Grounding Paradigm for Large Language Models in Recommendation Systems.
\textit{arXiv preprint arXiv:2308.08434}, 2023.

\bibitem{bao2024d3}
Keqin Bao, Jizhi Zhang, Yang Zhang, Xinyue Huo, Chong Chen, and Fuli Feng.
Decoding Matters: Addressing Amplification Bias and Homogeneity Issue in Recommendations for Large Language Models.
In \textit{Proceedings of EMNLP}, 2024.

\bibitem{chen2024sdpo}
Yuxin Chen, Junfei Tan, An Zhang, Zhengyi Yang, Leheng Sheng, Enzhi Zhang, Xiang Wang, and Tat-Seng Chua.
On Softmax Direct Preference Optimization for Recommendation.
In \textit{Advances in Neural Information Processing Systems 37 (NeurIPS 2024)}, 2024.

\bibitem{liang2026resid}
Yu Liang, Zhongjin Zhang, Yuxuan Zhu, Kerui Zhang, Zhiluohan Guo, Wenhang Zhou, Zonqi Yang, Kangle Wu, Yabo Ni, Anxiang Zeng, Cong Fu, Jianxin Wang, and Jiazhi Xia.
Rethinking Generative Recommender Tokenizer: Recsys-Native Encoding and Semantic Quantization Beyond LLMs.
\textit{arXiv preprint arXiv:2602.02338}, 2026.

\bibitem{fang2025hidvae}
Dengzhao Fang, Jingtong Gao, Chengcheng Zhu, Yu Li, Xiangyu Zhao, and Yi Chang.
HiD‑VAE: Interpretable Generative Recommendation via Hierarchical and Disentangled Semantic IDs.
\textit{arXiv preprint arXiv:2508.04618}, 2025.

\bibitem{wang2024letter}
Wenjie Wang, Honghui Bao, Xinyu Lin, Jizhi Zhang, Yongqi Li, Fuli Feng, See-Kiong Ng, and Tat-Seng Chua.
Learnable Item Tokenization for Generative Recommendation.
In \textit{Proceedings of the 33rd ACM International Conference on Information and Knowledge Management (CIKM 2024)}, 2024.

\bibitem{fang2026prism}
Dengzhao Fang, Jingtong Gao, Yu Li, Xiangyu Zhao, and Yi Chang.
PRISM: Purified Representation and Integrated Semantic Modeling for Generative Sequential Recommendation.
\textit{arXiv preprint arXiv:2601.16556}, 2026.

\end{thebibliography}


\end{document}
